\section{Síntesis y ampliación}

Esta clase amplía la seguridad infantil en línea: deja de tratarla como un problema aislado de clasificación de contenido y la presenta como un frontier system completo. Lo más importante no es el nombre de un producto, sino los siguientes principios, aplicables a otros contextos:

\begin{enumerate}
  \item \textbf{Prioridad al resultado de extremo a extremo}: signal, platform decision, lawful report y victim protection son etapas distintas y requieren un handoff claro entre ellas;
  \item \textbf{Separación en capas de lo known y lo unknown}: hash matching procesa a alta velocidad los objetos ya conocidos y predictive detection amplía la cobertura de candidatos desconocidos; ambos requieren policy y revisión;
  \item \textbf{Diseño conjunto del modelo y la cola}: precision, recall, calibration, threshold, review capacity y moderator wellbeing no pueden validarse por separado;
  \item \textbf{La gobernanza de datos es una capacidad central}: trusted source, provenance, segregation, retention, label protocol y audit determinan si el modelo puede desplegarse;
  \item \textbf{Safety by Design abarca todo el ciclo de vida}: development, deployment, maintenance, red teaming y transparency tienen cada uno un owner definido;
  \item \textbf{La privacidad es un objetivo de arquitectura}: E2EE, metadata, endpoint y network signal tienen cada uno sus propios límites de derechos y de seguridad; deben minimizarse y admitir apelación;
  \item \textbf{Las señales más tempranas exigen acciones más mesuradas}: conversation y graph analysis aportan alertas tempranas, pero las acciones deben ser escalonadas, reversibles y conservar la human review;
  \item \textbf{Los roles del ecosistema no son intercambiables}: plataformas, proveedores de tecnología, organismos receptores de reportes, investigadores y servicios de apoyo forman un ciclo cerrado de protección mediante contract y audit.
\end{enumerate}

\begin{importantbox}{Tarea de seguridad de plataformas: diseñar un ciclo cerrado verificable}
Elige un producto que permita a los usuarios subir contenido o enviar mensajes privados y entrega un diseño de una página: dibuja la máquina de estados de detección, revisión, actuación y reporte; define los campos de known-match y de predictive signal; establece las restricciones de precision/recall y de reviewer-capacity; redacta las políticas de privacy, retention, appeal, incident y provider-outage. No se permite dibujar solo un recuadro de «AI moderation»: debe indicarse el owner y la evidencia de cada handoff.
\end{importantbox}

\subsection{Lecturas adicionales}

{\small
\begin{itemize}[nosep,leftmargin=*]
  \item \href{https://www.thorn.org/solutions/for-platforms/}{\textbf{Thorn platform solutions}}: punto de acceso oficial a Safer detection, Match y Predict.
  \item \href{https://info.thorn.org/hubfs/thorn-safety-by-design-for-generative-AI.pdf}{\textbf{Thorn Safety by Design for Generative AI}}: documento técnico sobre principios para el ciclo de vida de la IA generativa.
  \item \href{https://www.missingkids.org/gethelpnow/cybertipline/cybertiplinedata}{\textbf{NCMEC CyberTipline data}}: criterios de los informes anuales y descripción de tendencias.
  \item \href{https://www.missingkids.org/netsmartz/home}{\textbf{NCMEC NetSmartz}}: recursos de educación preventiva dirigidos a familias, escuelas y niños.
  \item \href{https://nvlpubs.nist.gov/nistpubs/ai/NIST.AI.600-1.pdf}{\textbf{NIST AI 600-1}}: Generative AI Profile.
  \item \href{https://nvlpubs.nist.gov/nistpubs/ai/NIST.AI.100-4.pdf}{\textbf{NIST AI 100-4}}: synthetic-content risk technical approaches.
\end{itemize}
}

\end{document}
